\subsection*{AI Use Statement}

In this work, we used generative AI tools to assist with writing, information retrieval and discovery, and research ideation and execution.

For writing, generative AI was used to assist with drafting, polishing, and checking the formatting of the manuscript. For retrieval and discovery, generative AI was used to summarize literature and identify related work and relevant sources. For research ideation and execution, generative AI was used to assist with code understanding, simple script generation, code review, debugging, and identifying potential technical issues.

All AI-assisted outputs were reviewed by the authors. Suggested literature and references were checked against the original sources, AI-assisted code and debugging suggestions were manually inspected and validated through execution and testing, and AI-assisted manuscript text was reviewed and revised for technical correctness and consistency with the reported results.

We take responsibility for the final content of this work, including all text, claims, code, and other artifacts produced with the aid of generative AI.

% \subsection*{Ethics statement}

% (This section is \textbf{recommended} and does not count toward the page limit.)

% If authors feel that their paper submission raises questions regarding the Code
% of Ethics, they are encouraged to include a paragraph of Ethics Statement (at
% the end of the main text before references) to address potential concerns where
% appropriate. Topics include, but are not limited to, studies that involve human
% subjects, practices to data set releases, potentially harmful insights,
% methodologies and applications, potential conflicts of interest and sponsorship,
% discrimination/bias/fairness concerns, privacy and security issues, legal
% compliance, and research integrity issues (e.g., IRB, documentation, research
% ethics). This statement should not be more than 1 page.


\section*{Reproducibility Statement}

We make extensive efforts to facilitate the reproducibility of SpecSys.  Additional details required to reproduce our experiments are provided in the appendix. In particular, \autoref{app:hyperparameters} reports the agent-trace collection and replay procedure, inference-engine and generation configurations, speculative decoding settings, scheduling parameters, and software environments. \autoref{prompt_parser} provides concrete examples and implementation details of the feature extraction and remaining-step prediction procedures for different workflow types. We further report the workload and request characteristics, hardware serving capacity, and online trace statistics in~\autoref{request-and-resource-analysis} and~\autoref{online-trace-detail}. The profiling methodology and configuration-dependent thresholds used for selective SD are documented in~\autoref{sd-speedup}. Together, these descriptions specify the workloads, serving environments, system configurations, and implementation choices used throughout our evaluation. We have also released the SpecSys source code and the collected agent execution trajectories in an anonymous public repository at \url{https://anonymous.4open.science/r/SpecSys-Code-E03D/} to facilitate reproduction of our results.





% \subsubsection*{Acknowledgments}
% Use unnumbered third level headings for the acknowledgments. All
% acknowledgments, including those to funding agencies, go at the end of the paper.